\section*{Վարժություն 4. Ֆունկցիաներ (Արտապատկերումներ)}

% ============================================================
\subsection*{Խնդիր 1}
Եթե $f$-ը ամբողջական է, իսկ $g$-ն՝ ոչ (այսինքն՝ $g$-ն մասնակի ֆունկցիա է).
\begin{itemize}[nosep]
  \item $g \circ f$-ը կարող է լինել մասնակի. այն որոշված է $x$-ում միայն այն դեպքում, երբ $f(x) \in \operatorname{Dom}(g)$:
  \item $f \circ g$-ն որոշված է $x$-ում միայն այն դեպքում, երբ $g(x)$-ը որոշված է, բայց որտեղ $g(x)$-ը որոշված է, $f(g(x))$-ը նույնպես որոշված է, քանի որ $f$-ը ամբողջական է: Այսպիսով՝ $\operatorname{Dom}(f \circ g) = \operatorname{Dom}(g)$:
\end{itemize}

% ============================================================
\subsection*{Խնդիր 2}
\begin{enumerate}[label=(\alph*), nosep]
  \item $f(x)=x^2+3$:\quad $D_f = \mathbb{R}$,\; $R_f = [3,+\infty)$:
  \item $f(x)=\sqrt{x+2}$:\quad $D_f = [-2,+\infty)$,\; $R_f = [0,+\infty)$:
  \item $f(x)=\frac{1}{\sqrt{x-2}}$:\quad $D_f = (2,+\infty)$,\; $R_f = (0,+\infty)$:
  \item $f(x)=|x|$:\quad $D_f = \mathbb{R}$,\; $R_f = [0,+\infty)$:
  \item $f(x)=\frac{1}{x^2+2}$:\quad $D_f = \mathbb{R}$,\; $R_f = \bigl(0,\,\tfrac{1}{2}\bigr]$:
  \item $f(x)=\frac{1}{x^2-2}$:\quad $D_f = \mathbb{R}\setminus\{-\sqrt{2},\sqrt{2}\}$,\; $R_f = \bigl(-\infty,-\tfrac{1}{2}\bigr]\cup(0,+\infty)$:
\end{enumerate}

% ============================================================
\subsection*{Խնդիր 3}
\begin{enumerate}[label=(\alph*), nosep]
  \item $\rho = \{(x,\sqrt{x^2+4}) \mid x\in\mathbb{R}\}$: \textbf{Այո}, սա ֆունկցիա է: Յուրաքանչյուր $x$ արտապատկերվում է եզակի $\sqrt{x^2+4}$ արժեքի:
  \item $\rho = \{(x,5) \mid x\in\mathbb{R}\}$: \textbf{Այո}, սա (հաստատուն) ֆունկցիա է:
  \item $\rho = \{(7,y) \mid y\in\mathbb{R}\}$: \textbf{Ոչ}, սա ֆունկցիա չէ: $7$ տարրը զուգորդվում է յուրաքանչյուր իրական թվի հետ՝ խախտելով միարժեքությունը:
\end{enumerate}

% ============================================================
\subsection*{Խնդիր 4}
\begin{enumerate}[label=(\alph*), nosep]
  \item $f(x)=x^2+1$, $g(x)=x+3$:
    \[
      (f\circ g)(x) = x^2+6x+10, \qquad (g\circ f)(x) = x^2+4:
    \]
  \item $f(x)=\sqrt{x^2+2}$, $g(x)=x^2+3$:
    \[
      (f\circ g)(x) = \sqrt{x^4+6x^2+11}, \qquad (g\circ f)(x) = x^2+5:
    \]
  \item $f(x)=\frac{1}{x}$, $g(x)=2x+3$:
    \[
      (f\circ g)(x) = \frac{1}{2x+3}\;(x\neq -\tfrac{3}{2}), \qquad (g\circ f)(x) = \frac{2+3x}{x}\;(x\neq 0):
    \]
\end{enumerate}

% ============================================================
\subsection*{Խնդիր 5}
\begin{enumerate}[label=(\alph*), nosep]
  \item $f(x)=\frac{x+4}{2}$:\quad $f^{-1}(x)=2x-4$:
  \item $f(x)=x^3$:\quad $f^{-1}(x)=\sqrt[3]{x}$:
  \item $f(x)=\frac{x-2}{x+3}$:\quad $f^{-1}(x)=\dfrac{2+3x}{1-x},\; x\neq 1$:
\end{enumerate}

% ============================================================
\subsection*{Խնդիր 6}
\begin{enumerate}[label=(\alph*), nosep]
  \item $f(x)=|x|$: ոչ ինյեկտիվ է, ոչ էլ սյուրյեկտիվ: ($f(1)=f(-1)$; արժեքների տիրույթը $=[0,\infty)\neq\mathbb{R}$:)
  \item $f(x)=x^2+4$: ոչ ինյեկտիվ է, ոչ էլ սյուրյեկտիվ: ($f(1)=f(-1)$; արժեքների տիրույթը $=[4,\infty)\neq\mathbb{R}$:)
  \item $f(x)=x^3+6$: \textbf{բիյեկտիվ է}: (Խիստ աճող է, հետևաբար՝ ինյեկտիվ; խորանարդ արմատը ապահովում է սյուրյեկտիվություն:)
  \item $f(x)=x+|x|$: ոչ ինյեկտիվ է, ոչ էլ սյուրյեկտիվ: ($f(x)=0$ բոլոր $x<0$; արժեքների տիրույթը $=[0,\infty)\neq\mathbb{R}$:)
  \item $f(x)=x(x-2)(x+2)=x^3-4x$: սյուրյեկտիվ է, բայց \textbf{ոչ} ինյեկտիվ: (Խորանարդը ձգտում է $\pm\infty$, ինչը տալիս է սյուրյեկտիվություն; $f(0)=f(2)=0$ ցույց է տալիս, որ ինյեկտիվ չէ:)
\end{enumerate}

% ============================================================
\subsection*{Խնդիր 7}
Միջուկը՝ $\operatorname{Ker}(f)=\{(x_1,x_2): f(x_1)=f(x_2)\}$: Համարժեքության դասերն են.
\begin{enumerate}[label=(\alph*), nosep]
  \item $f:\mathbb{Z}\to\mathbb{Z}$, $f(x)=|x|$:\; $\operatorname{Ker}(f)$-ը տրոհում է $\mathbb{Z}$-ը $\{0\}$ և $\{-n,\,n\}$ բազմությունների՝ յուրաքանչյուր $n\geq 1$-ի համար:
  \item $f:\mathbb{R}\to\mathbb{Z}$, $f(x)=\lfloor x\rfloor$:\; դասերը $[n,\,n+1)$ կիսաբաց միջակայքերն են յուրաքանչյուր $n\in\mathbb{Z}$-ի համար:
  \item $f:\mathbb{R}\to\mathbb{Z}$, $f(x)=\lceil x\rceil$:\; դասերը $(n-1,\,n]$ կիսաբաց միջակայքերն են յուրաքանչյուր $n\in\mathbb{Z}$-ի համար:
\end{enumerate}

% ============================================================
\subsection*{Խնդիր 8}
Պետք է ցույց տանք, որ $(a)\Leftrightarrow(b)$, որտեղ (a)-ն $f(x)=f(y)\Rightarrow x=y$ պնդումն է, իսկ (b)-ն՝ $x\neq y\Rightarrow f(x)\neq f(y)$:

(b) պնդումը (a)-ի հակադիր (contrapositive) պնդումն է: Պայմանական պնդումը և դրա հակադիրը տրամաբանորեն համարժեք են, ուստի (a)$\Leftrightarrow$(b): \qed

% ============================================================
\subsection*{Խնդիր 9}
(Հղում. Լեմմաներ 1, 2, 3, 5 ֆունկցիաների մասին դասախոսության նշումներից:)

\textbf{Լեմմա 1} (Ինյեկցիաների համադրույթը ինյեկտիվ է). Եթե $f,g$ ինյեկտիվ են, ենթադրենք $(g\circ f)(x_1)=(g\circ f)(x_2)$: Այդ դեպքում՝ $g(f(x_1))=g(f(x_2))$; $g$ ինյեկտիվ $\Rightarrow f(x_1)=f(x_2)$; $f$ ինյեկտիվ $\Rightarrow x_1=x_2$: \qed

\textbf{Լեմմա 2} (Սյուրյեկցիաների համադրույթը սյուրյեկտիվ է). Եթե $f:A\to B$ և $g:B\to C$ սյուրյեկտիվ են, տրված $c\in C$-ի համար $\exists\, b\in B$, որ $g(b)=c$, և $\exists\, a\in A$, որ $f(a)=b$, ուստի $(g\circ f)(a)=c$: \qed

\textbf{Լեմմա 3} (Բիյեկցիաների համադրույթը բիյեկտիվ է). Հետևում է Լեմմա 1-ից և 2-ից:

\textbf{Լեմմա 5} (Եթե $g\circ f$-ը ինյեկտիվ է, ապա $f$-ը ինյեկտիվ է). Ենթադրենք $f(x_1)=f(x_2)$: Ապա $g(f(x_1))=g(f(x_2))$, այսինքն՝ $(g\circ f)(x_1)=(g\circ f)(x_2)$; քանի որ $g\circ f$-ը ինյեկտիվ է, $x_1=x_2$: \qed

% ============================================================
\subsection*{Խնդիր 10}
$\operatorname{Ker}(f) = \{(a_1,a_2)\in A\times A : f(a_1)=f(a_2)\}$-ը համարժեքության հարաբերություն է.
\begin{itemize}[nosep]
  \item \emph{Ռեֆլեքսիվ.} $f(a)=f(a)$, ուստի $a\sim a$:
  \item \emph{Սիմետրիկ.} $f(a)=f(b)\Rightarrow f(b)=f(a)$:
  \item \emph{Տրանզիտիվ.} $f(a)=f(b)$ և $f(b)=f(c)\Rightarrow f(a)=f(c)$: \qed
\end{itemize}

% ============================================================
\subsection*{Խնդիր 11}
\begin{enumerate}[label=(\alph*), nosep]
  \item $f:\mathbb{Z}\to\mathbb{N}$, $f(x)=|x|$:\; համարժեքության դասերն են $\{0\}$ և $\{-n,n\}$, որտեղ $n\geq 1$:
  \item $f:\mathbb{R}\to\mathbb{Z}$, $f(x)=\lfloor x\rfloor$:\; համարժեքության դասերը $[n,\,n+1)$ միջակայքերն են յուրաքանչյուր $n\in\mathbb{Z}$-ի համար:
\end{enumerate}
(Նույնն է, ինչ Խնդիր 7(a) և 7(b):)

% ============================================================
\subsection*{Խնդիր 12}
$f:\mathbb{Q}^+\to\mathbb{Q}^+$, $f(x)=\frac{x}{2x+1}$:

\textbf{Ինյեկտիվ է.} $\frac{x_1}{2x_1+1}=\frac{x_2}{2x_2+1}$ $\Rightarrow$ խաչաձև բազմապատկենք $\Rightarrow$ $2x_1x_2+x_1=2x_1x_2+x_2$ $\Rightarrow$ $x_1=x_2$: \checkmark

\textbf{Սյուրյեկտիվ չէ.} $f(x)=\frac{1}{2+1/x}<\frac{1}{2}$ բոլոր $x>0$-ների համար, ուստի $y=\frac{1}{2}\in\mathbb{Q}^+$ չունի նախապատկեր: \checkmark

% ============================================================
\subsection*{Խնդիր 13}
$f:\mathbb{Z}_{15}\to\mathbb{Z}_{15}$, $f(x)\equiv 4x\pmod{15}$:

\textbf{Բիյեկտիվ է.} $\gcd(4,15)=1$, ուստի $4$-ով բազմապատկումը բիյեկցիա է $\mathbb{Z}_{15}$-ի վրա:

\textbf{Հակադարձ.} Մեզ պետք է $4^{-1}\pmod{15}$: Քանի որ $4\cdot 4=16\equiv 1\pmod{15}$, ապա $4^{-1}\equiv 4$: Այսպիսով՝ $f^{-1}(x)\equiv 4x\pmod{15}$ (ֆունկցիան իր հակադարձն է):

% ============================================================
\subsection*{Խնդիր 14}
$f:\mathbb{Z}_{15}\to\mathbb{Z}_{15}$, $f(x)\equiv 2x+4\pmod{15}$:

\textbf{Բիյեկտիվ է:} Քանի որ $\gcd(2,15)=1$, $2$-ով բազմապատկումը հակադարձելի է $\pmod{15}$-ով, ուստի $f$-ը բիյեկցիա է: Ստուգում. $f(0)=4,\;f(1)=6,\;f(2)=8,\;\ldots,\;f(13)=0,\;f(14)=2$ --- բոլոր $15$ մնացորդները հանդիպում են:

\textbf{Հակադարձ.} $2^{-1}\equiv 8\pmod{15}$ (քանի որ $2\cdot 8=16\equiv 1$): Ուստի $f^{-1}(y)\equiv 8(y-4)\equiv 8y+13\pmod{15}$:

% ============================================================
\subsection*{Խնդիր 15}
\textbf{Պնդում.} Եթե $A$-ն վերջավոր է, ապա $f:A\to A$-ն բիյեկտիվ է այն և միայն այն դեպքում, երբ այն ինյեկտիվ է:

\emph{Ապացույցի ուրվագիծ.} ($\Rightarrow$) Բիյեկտիվ $\Rightarrow$ ինյեկտիվ ակնհայտորեն: ($\Leftarrow$) Եթե $f$-ը ինյեկտիվ է, ապա $|f(A)|=|A|$: Քանի որ $f(A)\subseteq A$ և $A$-ն վերջավոր է, $|f(A)|=|A|$ պայմանից հետևում է $f(A)=A$, ուստի $f$-ը սյուրյեկտիվ է, հետևաբար՝ բիյեկտիվ: \qed

% ============================================================
\subsection*{Խնդիր 16}
\textbf{Պնդում.} Եթե $A$-ն վերջավոր է, ապա $f:A\to A$-ն բիյեկտիվ է այն և միայն այն դեպքում, երբ այն սյուրյեկտիվ է:

\emph{Ապացույցի ուրվագիծ.} ($\Rightarrow$) Ակնհայտ է: ($\Leftarrow$) Եթե $f$-ը սյուրյեկտիվ է, ապա $f(A)=A$, ուստի $|f(A)|=|A|$: Ըստ Դիրիխլեի սկզբունքի (pigeonhole principle), վերջավոր բազմությունից իր վրա սյուրյեկցիան պետք է լինի ինյեկտիվ (հակառակ դեպքում երկու մուտքային արժեքներ կունենան նույն ելքային արժեքը՝ թողնելով առավելագույնը $|A|-1$ տարբեր ելքային արժեքներ, ինչը հակասում է սյուրյեկտիվությանը): \qed

% ============================================================
\subsection*{Խնդիր 17}
\textbf{Հետևանք 1} (դասախոսությունից). Վերջավոր $A$ բազմության համար $f:A\to A$ ինյեկտիվ է $\Leftrightarrow$ սյուրյեկտիվ է $\Leftrightarrow$ բիյեկտիվ է:

Սա Խնդիր 15-ի և 16-ի անմիջական համադրությունն է: \qed

% ============================================================
\subsection*{Խնդիր 18}
(Հղում. Թեորեմներ 6 և 7 դասախոսության նշումներից:)

\textbf{Թեորեմ 6} ($f$-ը ունի ձախ հակադարձ այն և միայն այն դեպքում, երբ $f$-ը ինյեկտիվ է). ($\Leftarrow$) Եթե $f:A\to B$ ինյեկտիվ է, սահմանենք $g:B\to A$-ն այսպես՝ $g(y)=$ այն եզակի $x$-ը, որի համար $f(x)=y$, եթե $y\in\operatorname{Ran}(f)$, և $g(y)=a_0$ (ֆիքսված տարր) հակառակ դեպքում: Այդ դեպքում $g\circ f = \operatorname{id}_A$: ($\Rightarrow$) Եթե $g\circ f=\operatorname{id}_A$, ապա $f(x_1)=f(x_2)\Rightarrow g(f(x_1))=g(f(x_2))\Rightarrow x_1=x_2$:

\textbf{Թեորեմ 7} ($f$-ը ունի աջ հակադարձ այն և միայն այն դեպքում, երբ $f$-ը սյուրյեկտիվ է). ($\Leftarrow$) Եթե $f$-ը սյուրյեկտիվ է, յուրաքանչյուր $y\in B$-ի համար ընտրենք որևէ $x_y\in f^{-1}(\{y\})$ և սահմանենք $h(y)=x_y$: Այդ դեպքում $f\circ h=\operatorname{id}_B$: ($\Rightarrow$) Եթե $f\circ h=\operatorname{id}_B$, ապա ցանկացած $y\in B$-ի համար $f(h(y))=y$, ուստի $y\in\operatorname{Ran}(f)$:

% ============================================================
\subsection*{Խնդիր 19}
$f(A\cap B)=f(A)\cap f(B)$ տեղի ունի բոլոր $A,B$-ների համար այն և միայն այն դեպքում, երբ $f$-ը \textbf{ինյեկտիվ է}:

\emph{Ուրվագիծ.} ($\Rightarrow$) Եթե $f$-ը ինյեկտիվ չէ, վերցնենք $A=\{x_1\}$, $B=\{x_2\}$, որտեղ $x_1\neq x_2$, $f(x_1)=f(x_2)$: Այդ դեպքում $f(\varnothing)=\varnothing\neq\{f(x_1)\}=f(A)\cap f(B)$: ($\Leftarrow$) Եթե $f$-ը ինյեկտիվ է, ապա $f(A)\cap f(B)\subseteq f(A\cap B)$ ներդրումը տեղի ունի, որովհետև $f(a)=f(b)\Rightarrow a=b$: \qed

% ============================================================
\subsection*{Խնդիր 20}
$f(A\setminus B)=f(A)\setminus f(B)$ տեղի ունի բոլոր $A,B$-ների համար այն և միայն այն դեպքում, երբ $f$-ը \textbf{ինյեկտիվ է}:

\emph{Ուրվագիծ.} ($\Rightarrow$) Եթե $f$-ը ինյեկտիվ չէ, վերցնենք $a\neq b$, որոնց համար $f(a)=f(b)$, $A=\{a,b\}$, $B=\{b\}$: Այդ դեպքում $f(A\setminus B)=f(\{a\})=\{f(a)\}$, բայց $f(A)\setminus f(B)=\{f(a)\}\setminus\{f(b)\}=\varnothing$: ($\Leftarrow$) Եթե $f$-ը ինյեկտիվ է և $y\in f(A\setminus B)$, ապա $y=f(a)$ եզակի $a\in A\setminus B$-ի համար; ինյեկտիվությունը երաշխավորում է, որ ոչ մի $b\in B$ չի արտապատկերվում $y$-ին, ուստի $y\notin f(B)$, ինչից հետևում է $y\in f(A)\setminus f(B)$: Մյուս ներդրումը տեղի ունի ցանկացած ֆունկցիայի համար: \qed

% ============================================================
\subsection*{Խնդիր 21}
\textbf{Այո:} Հնարավոր է, որ $X_1\neq X_2$, բայց $f(X_1)=f(X_2)$:

Օրինակ՝ $f:\{1,2\}\to\{a\}$, որտեղ $f(1)=f(2)=a$: Վերցնենք $X_1=\{1\}$, $X_2=\{1,2\}$: Այդ դեպքում $X_1\neq X_2$, բայց $f(X_1)=\{a\}=f(X_2)$:

% ============================================================
\subsection*{Խնդիր 22}
\textbf{Այո:} Չհատվող ոչ դատարկ ենթաբազմությունները կարող են ունենալ հավասար պատկերներ, եթե $f$-ը ինյեկտիվ չէ:

Օրինակ՝ $A=\{1,2,3,4\}$, $B=\{a,b\}$, $f(1)=a$, $f(2)=b$, $f(3)=a$, $f(4)=b$: Վերցնենք $X_1=\{1,2\}$, $X_2=\{3,4\}$: Այդ դեպքում $X_1\cap X_2=\varnothing$ և $f(X_1)=\{a,b\}=f(X_2)$:

% ============================================================
\subsection*{Խնդիր 23}
\textbf{Այո:} Հնարավոր է, որ $Y_1\neq Y_2$, բայց $f^{-1}(Y_1)=f^{-1}(Y_2)$:

Օրինակ՝ $f:\{1\}\to\{a,b\}$, որտեղ $f(1)=a$: Թող $Y_1=\{a\}$ և $Y_2=\{a,b\}$: Այդ դեպքում $Y_1\neq Y_2$, բայց $f^{-1}(Y_1)=\{1\}=f^{-1}(Y_2)$, քանի որ $b$-ն չունի նախապատկեր:

% ============================================================
\subsection*{Խնդիր 24}
\textbf{Այո}, դա միշտ ճիշտ է: Եթե $Y_1\cap Y_2=\varnothing$, ապա $f^{-1}(Y_1)\cap f^{-1}(Y_2)=\varnothing$:

\emph{Ապացույց.} $f^{-1}(Y_1)\cap f^{-1}(Y_2)=f^{-1}(Y_1\cap Y_2)=f^{-1}(\varnothing)=\varnothing$՝ օգտագործելով այն փաստը, որ նախապատկերները պահպանում են հատումները: \qed

% ============================================================
\subsection*{Խնդիր 25}
\textbf{Պնդում.} $f(f^{-1}(Y))\subseteq Y$ ցանկացած $f:A\to B$ և $Y\subseteq B$-ի համար:

\emph{Ապացույց.} Դիցուք $y\in f(f^{-1}(Y))$: Այդ դեպքում $y=f(x)$ որևէ $x\in f^{-1}(Y)$-ի համար, ինչը նշանակում է՝ $f(x)\in Y$, ուստի $y\in Y$: \qed

\textbf{Հակաօրինակ հավասարության համար.} $A=\{1\}$, $B=\{a,b\}$, $f(1)=a$, $Y=B$: Այդ դեպքում $f(f^{-1}(Y))=\{a\}\subsetneq\{a,b\}=Y$:

Հավասարությունը տեղի ունի բոլոր $Y\subseteq B$-ների համար այն և միայն այն դեպքում, երբ $f$-ը սյուրյեկտիվ է: